\chapter{Rings and Modules of Fractions}\label{chapter3}

The formation of rings of fractions and the associated process of localization are perhaps the most important technical tools in commutative algebra. They correspond in the algebro-geometric picture to concentrating attention on an open set or near a point, and the importance of these notions should be self-evident. This chapter gives the definitions and simple properties of the formation of fractions.

The procedure by which one constructs the rational field $\Q$ from the ring of integers $\Z$ (and embeds $\Z$ in $\Q$) extends easily to any integral domain $A$ and produces the field of fractions of $A$. The construction consists in taking all ordered pairs $(a, s)$ where $a, s \in A$ and $s \neq 0$, and setting up an equivalence relation between such pairs:
\[
    (a, s) \sim (b, t) \iff at - bs = 0.
\]
This works only if $A$ is an integral domain, because the verification that the relation is transitive involves canceling, i.e.\ the fact that $A$ has no zero-divisor $\neq 0$. However, it can be generalized as follows:

Let $A$ be any ring. A \emph{multiplicatively closed subset}\index{Multiplicatively closed set} of $A$ is a subset $S$ of $A$ such that $1 \in S$ and $S$ is closed under multiplication: in other words $S$ is a subsemigroup of the multiplicative semigroup of $A$. Define a relation $\sim$ on $A \times S$ as follows:
\[
    (a, s) \sim (b, t) \iff (at - bs)u = 0 \text{ for some } u \in S.
\]
Clearly this relation is reflexive and symmetric. To show that it is transitive, suppose $(a, s) \sim (b, t)$ and $(b, t) \sim (c, u)$. Then there exist $v, w \in S$ such that $(at - bs)v = 0$ and $(bu - ct)w = 0$. Eliminate $b$ from these two equations and we have $(au - cs)tvw = 0$. Since $S$ is closed under multiplication, we have $tvw \in S$, hence $(a, s) \sim (c, u)$. Thus we have an equivalence relation.

Let $a/s$ denote the equivalence class of $(a, s)$, and let $S^{-1}A$ denote the set of equivalence classes. We put a ring structure on $S^{-1}A$ by defining addition and multiplication of these ``fractions'' $a/s$ in the same way as in elementary algebra: that is,
\[
    (a/s) + (b/t) = (at + bs)/st, \qquad (a/s)(b/t) = ab/st.
\]

\begin{exercise*}
    Verify that these definitions are independent of the choices of representatives $(a, s)$ and $(b, t)$, and that $S^{-1}A$ satisfies the axioms of a commutative ring with identity.
\end{exercise*}

We also have a ring homomorphism $f \colon A \to S^{-1}A$ defined by $f(x) = x/1$. This is not in general injective.

\begin{remark}
If $A$ is an integral domain and $S = A \setminus \{0\}$, then $S^{-1}A$ is the field of fractions of $A$.
\end{remark}

The ring $S^{-1}A$ is called the \emph{ring of fractions}\index{Ring!of fractions}\index{Fractions, ring of} of $A$ with respect to $S$. It has a universal property:

\begin{proposition}\label{prop:universal-property}
Let $g \colon A \to B$ be a ring homomorphism such that $g(s)$ is a unit in $B$ for all $s \in S$. Then there exists a unique ring homomorphism $h \colon S^{-1}A \to B$ such that $g = h \circ f$.
\end{proposition}

\begin{proof}
i) \textit{Uniqueness.} If $h$ satisfies the conditions, then $h(a/1) = hf(a) = g(a)$ for all $a \in A$; hence, if $s \in S$,
\[
    h(1/s) = h((s/1)^{-1}) = h(s/1)^{-1} = g(s)^{-1}
\]
and therefore $h(a/s) = h(a/1) \cdot h(1/s) = g(a)g(s)^{-1}$, so that $h$ is uniquely determined by $g$.

ii) \textit{Existence.} Let $h(a/s) = g(a)g(s)^{-1}$. Then $h$ will clearly be a ring homomorphism provided that it is well-defined. Suppose then that $a/s = a'/s'$; then there exists $t \in S$ such that $(as' - a's)t = 0$, hence
\[
    (g(a)g(s') - g(a')g(s))g(t) = 0;
\]
now $g(t)$ is a unit in $B$, hence $g(a)g(s)^{-1} = g(a')g(s')^{-1}$.
\end{proof}

The ring $S^{-1}A$ and the homomorphism $f \colon A \to S^{-1}A$ have the following properties:
\begin{enumerate}[1)]
    \item $s \in S \implies f(s)$ is a unit in $S^{-1}A$;
    \item $f(a) = 0 \implies as = 0$ for some $s \in S$;
    \item Every element of $S^{-1}A$ is of the form $f(a)f(s)^{-1}$ for some $a \in A$ and some $s \in S$.
\end{enumerate}

Conversely, these three conditions determine the ring $S^{-1}A$ up to isomorphism. Precisely:

\begin{corollary}\label{cor:characterization-localization}
If $g \colon A \to B$ is a ring homomorphism such that
\begin{enumerate}[i)]
    \item $s \in S \implies g(s)$ is a unit in $B$;
    \item $g(a) = 0 \implies as = 0$ for some $s \in S$;
    \item Every element of $B$ is of the form $g(a)g(s)^{-1}$; then there is a unique isomorphism $h \colon S^{-1}A \to B$ such that $g = h \circ f$.
\end{enumerate}
\end{corollary}

\begin{proof}
By \eqref{prop:universal-property} we have to show that $h \colon S^{-1}A \to B$, defined by
\[
    h(a/s) = g(a)g(s)^{-1}
\]
(this definition uses i)) is an isomorphism. By iii), $h$ is surjective. To show $h$ is injective, look at the kernel of $h$: if $h(a/s) = 0$, then $g(a) = 0$, hence by ii) we have $at = 0$ for some $t \in S$, hence $(a, s) \sim (0, 1)$, i.e., $a/s = 0$ in $S^{-1}A$.
\end{proof}

\begin{examples}\label{ex:localization-examples}
\begin{enumerate}[(1)]
    \item \label{item 1-ex:localization-examples} Let $\mathfrak{p}$ be a prime ideal of $A$. Then $S = A - \mathfrak{p}$ is multiplicatively closed (in fact $A - \mathfrak{p}$ is multiplicatively closed $\iff \mathfrak{p}$ is prime). We write $A_{\mathfrak{p}}$ for $S^{-1}A$ in this case. The elements $a/s$ with $a \in \mathfrak{p}$ form an ideal $\mathfrak{m}$ in $A_{\mathfrak{p}}$. If $b/t \notin \mathfrak{m}$, then $b \notin \mathfrak{p}$, hence $b \in S$ and therefore $b/t$ is a unit in $A_{\mathfrak{p}}$. It follows that if $\mathfrak{a}$ is an ideal in $A_{\mathfrak{p}}$ and $\mathfrak{a} \not\subseteq \mathfrak{m}$, then $\mathfrak{a}$ contains a unit and is therefore the whole ring. Hence $\mathfrak{m}$ is the only maximal ideal in $A_{\mathfrak{p}}$; in other words, $A_{\mathfrak{p}}$ is a \emph{local ring}\index{Local ring}.

    The process of passing from $A$ to $A_{\mathfrak{p}}$ is called \emph{localization at $\mathfrak{p}$}\index{Localzation}.

    \item \label{item 2-ex:localization-examples} $S^{-1}A$ is the zero ring $\iff 0 \in S$.
    \item \label{item 3-ex:localization-examples} Let $f \in A$ and let $S = \{f^n\}_{n \geq 0}$. We write $A_f$ for $S^{-1}A$ in this case.
    \item \label{item 4-ex:localization-examples} Let $\mathfrak{a}$ be any ideal in $A$, and let $S = 1 + \mathfrak{a} = $ set of all $1 + x$ where $x \in \mathfrak{a}$. Clearly $S$ is multiplicatively closed.
    \item Special cases of \eqref{item 1-ex:localization-examples} and \eqref{item 3-ex:localization-examples}:
    \begin{enumerate}[i)]
        \item $A = \Z$, $\mathfrak{p} = (p)$, $p$ a prime number; $A_{\mathfrak{p}} = $ set of all rational numbers $m/n$ where $n$ is prime to $p$; if $f \in \Z$ and $f \neq 0$, then $A_f$ is the set of all rational numbers whose denominator is a power of $f$.

        \item $A = k[t_1, \ldots, t_n]$, where $k$ is a field and the $t_i$ are independent indeterminates, $\mathfrak{p}$ a prime ideal in $A$. Then $A_{\mathfrak{p}}$ is the ring of all rational functions $f/g$, where $g \notin \mathfrak{p}$. If $V$ is the variety defined by the ideal $\mathfrak{p}$, that is to say the set of all $x = (x_1, \ldots, x_n) \in k^n$ such that $f(x) = 0$ whenever $f \in \mathfrak{p}$, then (provided $k$ is infinite) $A_{\mathfrak{p}}$ can be identified with the ring of all rational functions on $k^n$ which are defined at almost all points of $V$; it is the \emph{local ring of $k^n$ along the variety $V$}. This is the prototype of the local rings which arise in algebraic geometry.
    \end{enumerate}
\end{enumerate}
\end{examples}

The construction of $S^{-1}A$ can be carried through with an $A$-module $M$ in place of the ring $A$. Define a relation $\sim$ on $M \times S$ as follows:
\[
    (m, s) \sim (m', s') \iff \exists\, t \in S \text{ such that } t(sm' - s'm) = 0.
\]
As before, this is an equivalence relation. Let $m/s$ denote the equivalence class of the pair $(m, s)$, let $S^{-1}M$ denote the set of such fractions, and make $S^{-1}M$ into an $S^{-1}A$-module with the obvious definitions of addition and scalar multiplication. As in Examples \eqref{item 1-ex:localization-examples} and \eqref{item 3-ex:localization-examples} above, we write $M_{\mathfrak{p}}$ instead of $S^{-1}M$ when $S = A - \mathfrak{p}$ ($\mathfrak{p}$ prime) and $M_f$ when $S = \{f^n\}_{n \geq 0}$.

Let $u \colon M \to N$ be an $A$-module homomorphism. Then it gives rise to an $S^{-1}A$-module homomorphism $S^{-1}u \colon S^{-1}M \to S^{-1}N$, namely $S^{-1}u$ maps $m/s$ to $u(m)/s$. We have $S^{-1}(v \circ u) = (S^{-1}v) \circ (S^{-1}u)$.

\begin{proposition}\label{prop:localization-exact}
The operation $S^{-1}$ is exact, i.e., if
\[
    M' \xrightarrow{f} M \xrightarrow{g} M''
\]
is exact at $M$, then
\[
    S^{-1}M' \xrightarrow{S^{-1}f} S^{-1}M \xrightarrow{S^{-1}g} S^{-1}M''
\]
is exact at $S^{-1}M$.
\end{proposition}

\begin{proof}
We have $g \circ f = 0$, hence $S^{-1}g \circ S^{-1}f = S^{-1}(0) = 0$, hence $\Image(S^{-1}f) \subseteq \Ker(S^{-1}g)$. To prove the reverse inclusion, let $m/s \in \Ker(S^{-1}g)$, then
\[
    g(m)/s = 0 \text{ in } S^{-1}M'',
\]
hence there exists $t \in S$ such that $tg(m) = 0$ in $M''$. But $tg(m) = g(tm)$ since $g$ is an $A$-module homomorphism, hence $tm \in \Ker(g) = \Image(f)$ and therefore $tm = f(m')$ for some $m' \in M'$. Hence in $S^{-1}M$ we have
\[
    m/s = f(m')/st = (S^{-1}f)(m'/st) \in \Image(S^{-1}f).
\]
Hence $\Ker(S^{-1}g) \subseteq \Image(S^{-1}f)$.
\end{proof}

In particular, it follows from \eqref{prop:localization-exact} that if $M'$ is a submodule of $M$, the mapping $S^{-1}M' \to S^{-1}M$ is injective and therefore $S^{-1}M'$ can be regarded as a submodule of $S^{-1}M$. With this convention,

\begin{corollary}\label{cor:localization-operations}
Formation of fractions commutes with formation of finite sums, finite intersections and quotients. Precisely, if $N, P$ are submodules of an $A$-module $M$, then
\begin{enumerate}[i)]
    \item $S^{-1}(N + P) = S^{-1}(N) + S^{-1}(P)$
    \item $S^{-1}(N \cap P) = S^{-1}(N) \cap S^{-1}(P)$
    \item the $S^{-1}A$-modules $S^{-1}(M/N)$ and $(S^{-1}M)/(S^{-1}N)$ are isomorphic,
\end{enumerate}
\end{corollary}

\begin{proof}
(i) follows readily from the definitions and (ii) is easy to verify: if $y/s = z/t$ ($y \in N$, $z \in P$, $s, t \in S$) then $u(ty - sz) = 0$ for some $u \in S$, hence $w = uty = usz \in N \cap P$ and therefore $y/s = w/stu \in S^{-1}(N \cap P)$. Consequently $S^{-1}N \cap S^{-1}P \subseteq S^{-1}(N \cap P)$, and the reverse inclusion is obvious.

(iii) Apply $S^{-1}$ to the exact sequence $0 \to N \to M \to M/N \to 0$.
\end{proof}

\begin{proposition}\label{prop:localization-tensor}
Let $M$ be an $A$-module. Then the $S^{-1}A$-modules $S^{-1}M$ and $S^{-1}A \otimes_A M$ are isomorphic; more precisely, there exists a unique isomorphism $f \colon S^{-1}A \otimes_A M \to S^{-1}M$ for which
\begin{equation}\label{eq:tensor-localization}
    f((a/s) \otimes m) = am/s \text{ for all } a \in A,\ m \in M,\ s \in S. \tag{1}
\end{equation}
\end{proposition}

\begin{proof}
The mapping $S^{-1}A \times M \to S^{-1}M$ defined by
\[
    (a/s, m) \mapsto am/s
\]
is $A$-bilinear, and therefore by the universal property \eqref{prop:universal-property} of the tensor product induces an $A$-homomorphism
\[
    f \colon S^{-1}A \otimes_A M \to S^{-1}M
\]
satisfying \eqref{eq:tensor-localization}. Clearly $f$ is surjective, and is uniquely defined by \eqref{eq:tensor-localization}.

Let $\sum_i (a_i/s_i) \otimes m_i$ be any element of $S^{-1}A \otimes M$. If $s = \prod_i s_i \in S$, $t_i = \prod_{j \neq i} s_j$, we have
\[
    \sum_i \frac{a_i}{s_i} \otimes m_i = \sum_i \frac{a_i t_i}{s} \otimes m_i = \sum_i \frac{1}{s} \otimes a_i t_i m_i = \frac{1}{s} \otimes \sum_i a_i t_i m_i.
\]
so that every element of $S^{-1}A \otimes M$ is of the form $(1/s) \otimes m$. Suppose that $f((1/s) \otimes m) = 0$. Then $m/s = 0$, hence $tm = 0$ for some $t \in S$, and therefore
\[
    \frac{1}{s} \otimes m = \frac{t}{st} \otimes m = \frac{1}{st} \otimes tm = \frac{1}{st} \otimes 0 = 0.
\]
Hence $f$ is injective and therefore an isomorphism.
\end{proof}

\begin{corollary}\label{cor:flat-localization}
$S^{-1}A$ is a flat $A$-module.
\end{corollary}

\begin{proof}
    \eqref{prop:localization-exact} and\eqref{prop:localization-tensor}.
\end{proof}

\begin{proposition}\label{prop:localization-tensor-product}
If $M, N$ are $A$-modules, there is a unique isomorphism of $S^{-1}A$-modules $f \colon S^{-1}M \otimes_{S^{-1}A} S^{-1}N \to S^{-1}(M \otimes_A N)$ such that
\[
    f((m/s) \otimes (n/t)) = (m \otimes n)/st.
\]
\end{proposition}

\begin{proof}
    Use \eqref{prop:localization-tensor} and the canonical isomorphisms of Chapter \ref{chapter2}.
\end{proof}

In particular, if $\mathfrak{p}$ is any prime ideal, then
\[
    M_{\mathfrak{p}} \otimes_{A_{\mathfrak{p}}} N_{\mathfrak{p}} \cong (M \otimes_A N)_{\mathfrak{p}}
\]
as $A_{\mathfrak{p}}$-modules.

\section*{Local Properties}
\addcontentsline{toc}{section}{Local Properties}
A property $P$ of a ring $A$ (or of an $A$-module $M$) is said to be a \emph{local property} if the following is true:
\[
    A \text{ (or } M) \text{ has } P \iff A_{\mathfrak{p}} \text{ (or } M_{\mathfrak{p}}) \text{ has } P, \text{ for each prime ideal } \mathfrak{p} \text{ of } A.
\]
The following propositions give examples of local properties:

\begin{proposition}\label{prop:zero-local}
Let $M$ be an $A$-module. Then the following are equivalent:
\begin{enumerate}[i)]
    \item $M = 0$;
    \item $M_{\mathfrak{p}} = 0$ for all prime ideals $\mathfrak{p}$ of $A$;
    \item $M_{\mathfrak{m}} = 0$ for all maximal ideals $\mathfrak{m}$ of $A$.
\end{enumerate}
\end{proposition}

\begin{proof}
Clearly (i) $\implies$ (ii) $\implies$ (iii). Suppose (iii) satisfied and $M \neq 0$. Let $x$ be a non-zero element of $M$, and let $\mathfrak{a} = \Ann(x)$; $\mathfrak{a}$ is an ideal $\neq (1)$, hence is contained in a maximal ideal $\mathfrak{m}$ by \eqref{cor:ideal-in-maximal}. Consider $x/1 \in M_{\mathfrak{m}}$. Since $M_{\mathfrak{m}} = 0$ we have $x/1 = 0$, hence $x$ is killed by some element of $A \setminus \mathfrak{m}$; but this is impossible since $\Ann(x) \subseteq \mathfrak{m}$.
\end{proof}

\begin{proposition}\label{prop:injective-surjective-local}
Let $\phi \colon M \to N$ be an $A$-module homomorphism. Then the following are equivalent:
\begin{enumerate}[i)]
    \item $\phi$ is injective;
    \item $\phi_{\mathfrak{p}} \colon M_{\mathfrak{p}} \to N_{\mathfrak{p}}$ is injective for each prime ideal $\mathfrak{p}$;
    \item $\phi_{\mathfrak{m}} \colon M_{\mathfrak{m}} \to N_{\mathfrak{m}}$ is injective for each maximal ideal $\mathfrak{m}$.
\end{enumerate}
Similarly with ``injective'' replaced by ``surjective'' throughout.
\end{proposition}

\begin{proof}
(i) $\implies$ (ii). $0 \to M \to N$ is exact, hence $0 \to M_{\mathfrak{p}} \to N_{\mathfrak{p}}$ is exact, i.e.\ $\phi_{\mathfrak{p}}$ is injective.

(ii) $\implies$ (iii) because a maximal ideal is prime.

(iii) $\implies$ (i). Let $M' = \Ker(\phi)$, then the sequence $0 \to M' \to M \to N$ is exact, hence $0 \to M'_{\mathfrak{m}} \to M_{\mathfrak{m}} \to N_{\mathfrak{m}}$ is exact by \eqref{prop:localization-exact} and therefore $M'_{\mathfrak{m}} \cong \Ker(\phi_{\mathfrak{m}}) = 0$ since $\phi_{\mathfrak{m}}$ is injective. Hence $M' = 0$ by \eqref{prop:zero-local}, hence $\phi$ is injective.

For the other part of the proposition, just reverse all the arrows.
\end{proof}

Flatness is a local property:

\begin{proposition}\label{prop:flat-local}
For any $A$-module $M$, the following statements are equivalent:
\begin{enumerate}[i)]
    \item $M$ is a flat $A$-module;
    \item $M_{\mathfrak{p}}$ is a flat $A_{\mathfrak{p}}$-module for each prime ideal $\mathfrak{p}$;
    \item $M_{\mathfrak{m}}$ is a flat $A_{\mathfrak{m}}$-module for each maximal ideal $\mathfrak{m}$.
\end{enumerate}
\end{proposition}

\begin{proof}
(i) $\implies$ (ii) by \eqref{prop:localization-tensor} and \eqref{ex:flat-base-change}.

(ii) $\implies$ (iii) O.K.

(iii) $\implies$ (i). If $N \to P$ is a homomorphism of $A$-modules, and $\mathfrak{m}$ is any maximal ideal of $A$, then
\begin{align*}
    N \to P \text{ injective } &\implies N_{\mathfrak{m}} \to P_{\mathfrak{m}} \text{ injective, by \eqref{prop:injective-surjective-local}} \\
    &\implies N_{\mathfrak{m}} \otimes_{A_{\mathfrak{m}}} M_{\mathfrak{m}} \to P_{\mathfrak{m}} \otimes_{A_{\mathfrak{m}}} M_{\mathfrak{m}} \text{ injective, by \eqref{prop:flat-equivalent}} \\
    &\implies (N \otimes_A M)_{\mathfrak{m}} \to (P \otimes_A M)_{\mathfrak{m}} \text{ injective, by \eqref{prop:localization-tensor-product}} \\
    &\implies N \otimes_A M \to P \otimes_A M \text{ injective, by \eqref{prop:injective-surjective-local}.}
\end{align*}
Hence $M$ is flat by \eqref{prop:flat-equivalent}.
\end{proof}

\section*{Extended and Contracted Ideals in Rings of Fractions}
\addcontentsline{toc}{section}{Extended and Contracted Ideals in Rings of Fractions}

Let $A$ be a ring, $S$ a multiplicatively closed subset of $A$ and $f \colon A \to S^{-1}A$ the natural homomorphism, defined by $f(a) = a/1$. Let $C$ be the set of contracted ideals in $A$, and let $E$ be the set of extended ideals in $S^{-1}A$ (cf.\ \eqref{prop:finitely-generated-extension}). If $\mathfrak{a}$ is an ideal in $A$, its extension $\mathfrak{a}^e$ in $S^{-1}A$ is $S^{-1}\mathfrak{a}$ (for any $y \in \mathfrak{a}^e$ is of the form $\sum a_i/s_i$, where $a_i \in \mathfrak{a}$ and $s_i \in S$; bring this fraction to a common denominator).

\begin{proposition}\label{prop:extension-ideals}
\begin{enumerate}[i)]
    \item \label{item 1:prop:extension-ideals}Every ideal in $S^{-1}A$ is an extended ideal.
    \item \label{item 2:prop:extension-ideals}If $\mathfrak{a}$ is an ideal in $A$, then $\mathfrak{a}^{ec} = \bigcup_{s \in S} (\mathfrak{a} : s)$. Hence $\mathfrak{a}^e = (1)$ if and only if $\mathfrak{a}$ meets $S$.
    \item \label{item 3:prop:extension-ideals}$\mathfrak{a} \in C \iff$ no element of $S$ is a zero-divisor in $A/\mathfrak{a}$.
    \item \label{item 4:prop:extension-ideals}The prime ideals of $S^{-1}A$ are in one-to-one correspondence ($\mathfrak{p} \mapsto S^{-1}\mathfrak{p}$) with the prime ideals of $A$ which don't meet $S$.
    \item \label{item 5:prop:extension-ideals}The operation $S^{-1}$ commutes with formation of finite sums, products, intersections and radicals.
\end{enumerate}
\end{proposition}

\begin{proof}
    \begin{enumerate}[i)]
        \item Let $\mathfrak{b}$ be an ideal in $S^{-1}A$, and let $x/s \in \mathfrak{b}$. Then $x/1 \in \mathfrak{b}$, hence $x \in \mathfrak{b}^c$ and therefore $x/s \in \mathfrak{b}^{ec}$. Since $\mathfrak{b} \supseteq \mathfrak{b}^{ce}$ in any case \eqref{prop:extension-contraction-properties}, it follows that $\mathfrak{b} = \mathfrak{b}^{ce}$.
        \item $x \in \mathfrak{a}^{ec} = (S^{-1}\mathfrak{a})^c \iff x/1 = a/s$ for some $a \in \mathfrak{a}$, $s \in S \iff (xs - a)t = 0$ for some $t \in S \iff xst \in \mathfrak{a} \iff x \in \bigcup_{s \in S} (\mathfrak{a} : s)$.
        \item $\mathfrak{a} \in C \iff \mathfrak{a}^{ec} \subseteq \mathfrak{a} \iff (sx \in \mathfrak{a}$ for some $s \in S \implies x \in \mathfrak{a}) \iff$ no $s \in S$ is a zero-divisor in $A/\mathfrak{a}$.
        \item If $\mathfrak{q}$ is a prime ideal in $S^{-1}A$, then $\mathfrak{q}^c$ is a prime ideal in $A$ (this much is true for any ring homomorphism). Conversely, if $\mathfrak{p}$ is a prime ideal in $A$, then $A/\mathfrak{p}$ is an integral domain; if $S$ is the image of $S$ in $A/\mathfrak{p}$, we have $S^{-1}A/S^{-1}\mathfrak{p} \cong S^{-1}(A/\mathfrak{p})$ which is either $0$ or else is contained in the field of fractions of $A/\mathfrak{p}$ and is therefore an integral domain, and therefore $S^{-1}\mathfrak{p}$ is either prime or is the unit ideal; by (i) the latter possibility occurs if and only if $\mathfrak{p}$ meets $S$.
        \item For sums and products, this follows from \eqref{ex:extension-contraction-operations}; for intersections, from \eqref{cor:localization-operations}. As to radicals, we have $S^{-1}r(\mathfrak{a}) \subseteq r(S^{-1}\mathfrak{a})$ from \eqref{ex:extension-contraction-operations}, and the proof of the reverse inclusion is a routine verification which we leave to the reader. \qedhere
    \end{enumerate}
\end{proof}

\begin{remark}
\begin{enumerate}[1)]
    \item If $\mathfrak{a}, \mathfrak{b}$ are ideals of $A$, the formula
    \[
        S^{-1}(\mathfrak{a} : \mathfrak{b}) = (S^{-1}\mathfrak{a} : S^{-1}\mathfrak{b})
    \]
    is true provided the ideal $\mathfrak{b}$ is finitely generated: see \eqref{cor:localization-colon}.

    \item The proof in \eqref{prop:nilradical-intersection} that if $f \in A$ is not nilpotent there is a prime ideal of $A$ which does not contain $f$ can be expressed more concisely in the language of rings of fractions. Since the set $S = \{f^n\}_{n \geq 0}$ does not contain $0$, the ring $S^{-1}A = A_f$ is not the zero ring and therefore by \eqref{thm:maximal-ideal-exists} has a maximal ideal, whose contraction in $A$ is a prime ideal $\mathfrak{p}$ which does not meet $S$ by \eqref{prop:extension-ideals}; hence $f \notin \mathfrak{p}$.
\end{enumerate}
\end{remark}

\begin{corollary}\label{cor:localization-nilradical}
If $\mathfrak{R}$ is the nilradical of $A$, the nilradical of $S^{-1}A$ is $S^{-1}\mathfrak{R}$.
\end{corollary}

\begin{corollary}\label{cor:spec-localization}
If $\mathfrak{p}$ is a prime ideal of $A$, the prime ideals of the local ring $A_{\mathfrak{p}}$ are in one-to-one correspondence with the prime ideals of $A$ contained in $\mathfrak{p}$.
\end{corollary}

\begin{proof}
Take $S = A - \mathfrak{p}$ in \eqref{prop:extension-ideals} (iv).
\end{proof}

\begin{remark}
Thus the passage from $A$ to $A_{\mathfrak{p}}$ cuts out all prime ideals except those contained in $\mathfrak{p}$. In the other direction, the passage from $A$ to $A/\mathfrak{p}$ cuts out all prime ideals except those containing $\mathfrak{p}$. Hence if $\mathfrak{p}, \mathfrak{q}$ are prime ideals such that $\mathfrak{p} \supseteq \mathfrak{q}$, then by localizing with respect to $\mathfrak{p}$ and taking the quotient mod $\mathfrak{q}$ (in either order: these two operations commute, by \eqref{cor:localization-operations}), we restrict our attention to those prime ideals which lie between $\mathfrak{p}$ and $\mathfrak{q}$. In particular, if $\mathfrak{p} = \mathfrak{q}$ we end up with a field, called the \emph{residue field at $\mathfrak{p}$}, which can be obtained either as the field of fractions of the integral domain $A/\mathfrak{p}$ or as the residue field of the local ring $A_{\mathfrak{p}}$.
\end{remark}

\begin{proposition}\label{prop:localization-annihilator}
    Let $M$ be a finitely generated $A$-module, $S$ a multiplicatively closed subset of $A$. Then $S^{-1}(\Ann(M)) = \Ann(S^{-1}M)$.
\end{proposition}

\begin{proof}
    If this is true for two $A$-modules $M, N$, it is true for $M + N$:
    \begin{align*}
        S^{-1}(\Ann(M + N)) &= S^{-1}(\Ann(M) \cap \Ann(N)) && \text{by \eqref{ex:annihilator-properties}} \\
        &= S^{-1}(\Ann(M)) \cap S^{-1}(\Ann(N)) && \text{by \eqref{cor:localization-operations}} \\
        &= \Ann(S^{-1}M) \cap \Ann(S^{-1}N) && \text{by hypothesis} \\
        &= \Ann(S^{-1}M + S^{-1}N) = \Ann(S^{-1}(M + N)).
    \end{align*}
    Hence it is enough to prove \eqref{prop:localization-annihilator} for $M$ generated by a single element: then $M \cong A/\mathfrak{a}$ (as $A$-module), where $\mathfrak{a} = \Ann(M)$; $S^{-1}M \cong (S^{-1}A)/(S^{-1}\mathfrak{a})$ by \eqref{cor:localization-operations}, so that $\Ann(S^{-1}M) = S^{-1}\mathfrak{a} = S^{-1}(\Ann(M))$.
\end{proof}

\begin{corollary}\label{cor:localization-colon}
    If $N, P$ are submodules of an $A$-module $M$ and if $P$ is finitely generated, then $S^{-1}(N : P) = (S^{-1}N : S^{-1}P)$.
\end{corollary}

\begin{proof}
    $(N : P) = \Ann((N + P)/N)$ by \eqref{ex:annihilator-properties}; now apply \eqref{prop:localization-annihilator}.
\end{proof}

\begin{proposition}\label{prop:contracted-ideal}
Let $A \to B$ be a ring homomorphism and let $\mathfrak{p}$ be a prime ideal of $A$. Then $\mathfrak{p}$ is the contraction of a prime ideal of $B$ if and only if $\mathfrak{p}^{ec} = \mathfrak{p}$.
\end{proposition}

\begin{proof}
If $\mathfrak{p} = \mathfrak{q}^c$ then $\mathfrak{p}^{ec} = \mathfrak{p}$ by \eqref{prop:extension-contraction-properties}. Conversely, if $\mathfrak{p}^{ec} = \mathfrak{p}$, let $S$ be the image of $A - \mathfrak{p}$ in $B$. Then $\mathfrak{p}^e$ does not meet $S$, therefore by \eqref{prop:extension-ideals} its extension in $S^{-1}B$ is a proper ideal and hence is contained in a maximal ideal $\mathfrak{m}$ of $S^{-1}B$. If $\mathfrak{q}$ is the contraction of $\mathfrak{m}$ in $B$, then $\mathfrak{q}$ is prime, $\mathfrak{q} \supseteq \mathfrak{p}^e$ and $\mathfrak{q} \cap S = \varnothing$. Hence $\mathfrak{q}^c = \mathfrak{p}$.
\end{proof}

%=============================================================================
% Chapter 3 Exercises - With hints from original book (30 exercises)
%=============================================================================
\section*{Exercises}
\addcontentsline{toc}{section}{Exercises}

\begin{enumerate}
    \item \label{exer1-chapter3}Let $S$ be a multiplicatively closed subset of a ring $A$, and let $M$ be a finitely generated $A$-module. Prove that $S^{-1}M = 0$ if and only if there exists $s \in S$ such that $sM = 0$.

    \item \label{exer2-chapter3}Let $\mathfrak{a}$ be an ideal of a ring $A$, and let $S = 1 + \mathfrak{a}$. Show that $S^{-1}\mathfrak{a}$ is contained in the Jacobson radical of $S^{-1}A$.\\
    \textit{Hint:} Use this result and Nakayama's lemma to give a proof of \eqref{cor:cayley-hamilton} which does not depend on determinants. [If $M = \mathfrak{a}M$, then $S^{-1}M = (S^{-1}\mathfrak{a})(S^{-1}M)$, hence by Nakayama we have $S^{-1}M = 0$. Now use Exercise \ref{exer1-chapter3}.]

    \item \label{exer3-chapter3}Let $A$ be a ring, let $S$ and $T$ be two multiplicatively closed subsets of $A$, and let $U$ be the image of $T$ in $S^{-1}A$. Show that the rings $(ST)^{-1}A$ and $U^{-1}(S^{-1}A)$ are isomorphic.

    \item \label{exer4-chapter3}Let $f \colon A \to B$ be a homomorphism of rings and let $S$ be a multiplicatively closed subset of $A$. Let $T = f(S)$. Show that $S^{-1}B$ and $T^{-1}B$ are isomorphic as $S^{-1}A$-modules.

    \item \label{exer5-chapter3}Let $A$ be a ring. Suppose that, for each prime ideal $\mathfrak{p}$, the local ring $A_{\mathfrak{p}}$ has no nilpotent element $\neq 0$. Show that $A$ has no nilpotent element $\neq 0$. If each $A_{\mathfrak{p}}$ is an integral domain, is $A$ necessarily an integral domain?

    \item \label{exer6-chapter3}Let $A$ be a ring $\neq 0$ and let $\Sigma$ be the set of all multiplicatively closed subsets $S$ of $A$ such that $0 \notin S$. Show that $\Sigma$ has maximal elements, and that $S \in \Sigma$ is maximal if and only if $A \setminus S$ is a minimal prime ideal of $A$".

    \item \label{exer7-chapter3}A multiplicatively closed subset $S$ of a ring $A$ is said to be \emph{saturated}\index{Saturated} if $xy \in S \iff x \in S \text{ and } y \in S$. Prove that:
    \begin{enumerate}[i)]
        \item $S$ is saturated $\iff$ $A \setminus S$ is a union of prime ideals.
        \item If $S$ is any multiplicatively closed subset of $A$, there is a unique smallest saturated multiplicatively closed subset $\bar{S}$ containing $S$, and that $\bar{S}$ is the complement in $A$ of the union of the prime ideals which do not meet $S$. ($\bar{S}$ is called the \emph{saturation} of $S$.)
    \end{enumerate}
    If $S = 1 + \mathfrak{a}$, where $\mathfrak{a}$ is an ideal of $A$, find $\bar{S}$.

    \item \label{exer8-chapter3} Let $S, T$ be multiplicatively closed subsets of $A$, such that $S \subseteq T$. Let $\phi \colon S^{-1}A \to T^{-1}A$ be the homomorphism which maps each $a/s \in S^{-1}A$ to $a/s$ considered as an element of $T^{-1}A$. Show that the following statements are equivalent:
    \begin{enumerate}[i)]
        \item $\phi$ is bijective.
        \item For each $t \in T$, $t/1$ is a unit in $S^{-1}A$.
        \item For each $t \in T$ there exists $x \in A$ such that $xt \in S$.
        \item $T$ is contained in the saturation of $S$ (Exercise \ref{exer7-chapter3}).
        \item Every prime ideal which meets $T$ also meets $S$.
    \end{enumerate}

    \item \label{exer9-chapter3} The set $S_0$ of all non-zero-divisors in $A$ is a saturated multiplicatively closed subset of $A$. Hence the set $D$ of zero-divisors in $A$ is a union of prime ideals (see Chapter \ref{chapter1}, Exercise \ref{exer14-chapter1}). Show that every minimal prime ideal of $A$ is contained in $D$. \textit{Hint}: Use Exercise \ref{exer6-chapter3}.
    
    \quad\, The ring $S_0^{-1}A$ is called the \emph{total ring of fractions} of $A$. Prove that:
    \begin{enumerate}[i)]
        \item $S_0$ is the largest multiplicatively closed subset of $A$ for which the homomorphism $A \to S_0^{-1}A$ is injective.
        \item Every element in $S_0^{-1}A$ is either a zero-divisor or a unit.
        \item Every ring in which every non-unit is a zero-divisor is equal to its total ring of fractions (i.e., $A \to S_0^{-1}A$ is bijective).
    \end{enumerate}

    \item \label{exer10-chapter3} Let $A$ be a ring.
    \begin{enumerate}[i)]
        \item If $A$ is absolutely flat (Chapter \ref{chapter2}, Exercise \ref{exer27-chapter2}) and $S$ is any multiplicatively closed subset of $A$, then $S^{-1}A$ is absolutely flat.
        \item $A$ is absolutely flat $\iff$ $A_{\mathfrak{m}}$ is a field for each maximal ideal $\mathfrak{m}$.
    \end{enumerate}

    \item \label{exer11-chapter3} Let $A$ be a ring. Prove that the following are equivalent:
    \begin{enumerate}[i)]
        \item $A/\mathfrak{N}$ is absolutely flat ($\mathfrak{N}$ being the nilradical of $A$).
        \item Every prime ideal of $A$ is maximal.
        \item $\Spec(A)$ is a $T_1$-space (i.e., every subset consisting of a single point is closed).
        \item $\Spec(A)$ is Hausdorff.
    \end{enumerate}
    If these conditions are satisfied, show that $\Spec(A)$ is compact and totally disconnected (i.e., the only connected subsets of $\Spec(A)$ are those consisting of a single point).

    \item \label{exer12-chapter3} Let $A$ be an integral domain and $M$ an $A$-module. An element $x \in M$ is a \emph{torsion element}\index{Torsion, element} of $M$ if $\Ann(x) \neq 0$, that is if $x$ is killed by some non-zero element of $A$. Show that the torsion elements of $M$ form a submodule of $M$. This submodule is called the \emph{torsion submodule}\index{Torsion, element!submodule} of $M$ and is denoted by $T(M)$. If $T(M) = 0$, the module $M$ is said to be \emph{torsion-free}. Show that:
    \begin{enumerate}[i)]
        \item If $M$ is any $A$-module, then $M/T(M)$ is torsion-free.
        \item If $f \colon M \to N$ is a module homomorphism, then $f(T(M)) \subseteq T(N)$.
        \item If $0 \to M' \to M \to M''$ is an exact sequence, then the sequence $0 \to T(M') \to T(M) \to T(M'')$ is exact.
        \item If $M$ is any $A$-module, then $T(M)$ is the kernel of the mapping $x \mapsto 1 \otimes x$ of $M$ into $K \otimes_A M$, where $K$ is the field of fractions of $A$.
        
        \textit{Hint:} For iv), show that $K$ may be regarded as the direct limit of its submodules $Ag$ ($g \in K$); using Chapter \ref{chapter1}, Exercise \ref{exer15-chapter1} and Exercise \ref{exer20-chapter1}, show that if $1 \otimes x = 0$ in $K \otimes M$ then $1 \otimes x = 0$ in $Ag \otimes M$ for some $g \neq 0$. Deduce that $g^{-1}x = 0$.
    \end{enumerate}

    \item \label{exer13-chapter3} Let $S$ be a multiplicatively closed subset of an integral domain $A$. In the notation of Exercise \ref{exer12-chapter3}, show that $T(S^{-1}M) = S^{-1}(T(M))$. Deduce that the following are equivalent:
    \begin{enumerate}[i)]
        \item $M$ is torsion-free.
        \item $M_{\mathfrak{p}}$ is torsion-free for all prime ideals $\mathfrak{p}$.
        \item $M_{\mathfrak{m}}$ is torsion-free for all maximal ideals $\mathfrak{m}$.
    \end{enumerate}

    \item \label{exer14-chapter3} Let $M$ be an $A$-module and $\mathfrak{a}$ an ideal of $A$. Suppose that $M_{\mathfrak{m}} = 0$ for all maximal ideals $\mathfrak{m} \not\supseteq \mathfrak{a}$. Prove that $M = \mathfrak{a}M$.
    
    \textit{Hint:} Pass to the $A/\mathfrak{a}$-module $M/\mathfrak{a}M$ and use \eqref{prop:zero-local}.

    \item \label{exer15-chapter3}Let $A$ be a ring, and let $F$ be the $A$-module $A^n$. Show that every set of $n$ generators of $F$ is a basis of $F$.
    
    \textit{Hint:} Let $x_1, \ldots, x_n$ be a set of generators and $e_1, \ldots, e_n$ the canonical basis of $F$. Define $\phi \colon F \to F$ by $\phi(e_i) = x_i$. Then $\phi$ is surjective and we have to prove that it is an isomorphism. By \eqref{prop:injective-surjective-local} we may assume that $A$ is a local ring. Let $N$ be the kernel of $\phi$ and let $k = A/\mathfrak{m}$ be the residue field of $A$. Since $F$ is a flat $A$-module, the exact sequence $0 \to N \to F \xrightarrow{\phi} F \to 0$ gives an exact sequence $0 \longrightarrow k \otimes N \longrightarrow k \otimes F \xrightarrow{1 \otimes \phi} k \otimes F \longrightarrow 0$. Now $k \otimes F = k^n$ is an $n$-dimensional vector space over $k$; $1 \otimes \phi$ is surjective, hence bijective, hence $k \otimes N = 0$. Also $N$ is finitely generated, by Chapter \ref{chapter2}, Exercise \ref{exer12-chapter2}, hence $N = 0$ by Nakayama's lemma. Hence $\phi$ is an isomorphism.

    Deduce that every set of generators of $F$ has at least $n$ elements.

    \item \label{exer16-chapter3}Let $B$ be a flat $A$-algebra. Then the following conditions are equivalent:
    \begin{enumerate}[i)]
        \item $\mathfrak{a}^{ec} = \mathfrak{a}$ for all ideals $\mathfrak{a}$ of $A$.
        \item $\Spec(B) \to \Spec(A)$ is surjective.
        \item For every maximal ideal $\mathfrak{m}$ of $A$ we have $\mathfrak{m}^e \neq (1)$.
        \item If $M$ is any non-zero $A$-module, then $M_B \neq 0$.
        \item For every $A$-module $M$, the mapping $x \mapsto 1 \otimes x$ of $M$ into $M_B$ is injective.
    \end{enumerate}
    \textit{Hint:} For i) $\implies$ ii), use \eqref{prop:extension-ideals}. 
    
    ii) $\implies$ iii) is clear. 
    
    iii) $\implies$ iv): Let $x$ be a non-zero element of $M$ and let $M' = Ax$. Since $B$ is flat over $A$ it is enough to show that $M'_B \neq 0$. We have $M' \cong A/\mathfrak{a}$ for some ideal $\mathfrak{a} \neq (1)$, hence $M'_B \cong B/\mathfrak{a}^e$. Now $\mathfrak{a} \subseteq \mathfrak{m}$ for some maximal ideal $\mathfrak{m}$, hence $\mathfrak{a}^e \subseteq \mathfrak{m}^e \neq (1)$. Hence $M'_B \neq 0$. 
    
    iv) $\implies$ v): Let $M'$ be the kernel of $M \to M_B$. Since $B$ is flat over $A$, the sequence $0 \to M'_B \to M_B \to (M_B)_B$ is exact. But (Chapter \ref{chapter2}, Exercise \ref{exer13-chapter2}, with $N = M_B$) the mapping $M_B \to (M_B)_B$ is injective, hence $M'_B = 0$ and therefore $M' = 0$. 
    
    v) $\implies$ i): Take $M = A/\mathfrak{a}$.

    $B$ is said to be \emph{faithfully flat}\index{Flat!faithfully} over $A$.

    \item \label{exer17-chapter3}Let $A \xrightarrow{f} B \xrightarrow{g} C$ be ring homomorphisms. If $g \circ f$ is flat and $g$ is faithfully flat, then $f$ is flat.

    \item \label{exer18-chapter3}Let $f \colon A \to B$ be a flat homomorphism of rings, let $\mathfrak{q}$ be a prime ideal of $B$ and let $\mathfrak{p} = \mathfrak{q}^c$. Then $f^* \colon \Spec(B_{\mathfrak{q}}) \to \Spec(A_{\mathfrak{p}})$ is surjective.\\
    \textit{Hint:} For $B_{\mathfrak{p}}$ is flat over $A_{\mathfrak{p}}$ by \eqref{prop:flat-local}, and $B_{\mathfrak{q}}$ is a local ring of $B_{\mathfrak{p}}$, hence is flat over $B_{\mathfrak{p}}$. Hence $B_{\mathfrak{q}}$ is flat over $A_{\mathfrak{p}}$ and satisfies condition (3) of Exercise \ref{exer16-chapter3}.

    \item \label{exer19-chapter3}Let $A$ be a ring, $M$ an $A$-module. The \emph{support}\index{Support} of $M$ is defined to be the set $\Supp(M)$ of prime ideals $\mathfrak{p}$ of $A$ such that $M_{\mathfrak{p}} \neq 0$. Prove the following results:
    \begin{enumerate}[i)]
        \item $M \neq 0 \iff \Supp(M) \neq \varnothing$.
        \item $V(\mathfrak{a}) = \Supp(A/\mathfrak{a})$.
        \item If $0 \to M' \to M \to M'' \to 0$ is an exact sequence, then $\Supp(M) = \Supp(M') \cup \Supp(M'')$.
        \item If $M = \bigoplus M_i$, then $\Supp(M) = \bigcup \Supp(M_i)$.
        \item If $M$ is finitely generated, then $\Supp(M) = V(\Ann(M))$ (and is therefore a closed subset of $\Spec(A)$).
        \item If $M, N$ are finitely generated, then $\Supp(M \otimes_A N) = \Supp(M) \cap \Supp(N)$. [Use Chapter \ref{chapter2}, Exercise \ref{exer3-chapter2}.]
        \item If $M$ is finitely generated and $\mathfrak{a}$ is an ideal of $A$, then $\Supp(M/\mathfrak{a}M) = V(\mathfrak{a} + \Ann(M))$.
        \item If $f \colon A \to B$ is a ring homomorphism and $M$ is a finitely generated $A$-module, then $\Supp(B \otimes_A M) = f^{*-1}(\Supp(M))$.
    \end{enumerate}

    \item \label{exer20-chapter3}Let $f \colon A \to B$ be a ring homomorphism, $f^* \colon \Spec(B) \to \Spec(A)$ the associated mapping. Show that:
    \begin{enumerate}[i)]
        \item Every prime ideal of $A$ is a contracted ideal $\iff$ $f^*$ is surjective.
        \item Every prime ideal of $B$ is an extended ideal $\implies$ $f^*$ is injective.
    \end{enumerate}
    Is the converse of ii) true?

    \item \label{exer21-chapter3} \begin{enumerate}[i)]
        \item Let $A$ be a ring, $S$ a multiplicatively closed subset of $A$, and $\phi \colon A \to S^{-1}A$ the canonical homomorphism. Show that $\phi^* \colon \Spec(S^{-1}A) \to \Spec(A)$ is a homeomorphism of $\Spec(S^{-1}A)$ onto its image in $X = \Spec(A)$. Let this image be denoted by $S^{-1}X$. In particular, if $f \in A$, the image of $\Spec(A_f)$ in $X$ is the basic open set $X_f$ (Chapter \ref{chapter1}, Exercise \ref{exer17-chapter1}).
        \item Let $f \colon A \to B$ be a ring homomorphism. Let $X = \Spec(A)$ and $Y = \Spec(B)$, and let $f^* \colon Y \to X$ be the mapping associated with $f$. Identifying $\Spec(S^{-1}A)$ with its canonical image $S^{-1}X$ in $X$, and $\Spec(S^{-1}B)$ ($= \Spec(f(S)^{-1}B)$) with its canonical image $S^{-1}Y$ in $Y$, show that $S^{-1}f^* \colon \Spec(S^{-1}B) \to \Spec(S^{-1}A)$ is the restriction of $f^*$ to $S^{-1}Y$, and that $S^{-1}Y = f^{*-1}(S^{-1}X)$.
        \item Let $\mathfrak{a}$ be an ideal of $A$ and let $\mathfrak{b} = \mathfrak{a}^e$ be its extension in $B$. Let $\bar{f} \colon A/\mathfrak{a} \to B/\mathfrak{b}$ be the homomorphism induced by $f$. If $\Spec(A/\mathfrak{a})$ is identified with its canonical image $V(\mathfrak{a})$ in $X$, and $\Spec(B/\mathfrak{b})$ with its image $V(\mathfrak{b})$ in $Y$, show that $\bar{f}^*$ is the restriction of $f^*$ to $V(\mathfrak{b})$.
        \item Let $\mathfrak{p}$ be a prime ideal of $A$. Take $S = A \setminus \mathfrak{p}$ in ii) and then reduce mod $S^{-1}\mathfrak{p}$ as in iii). Deduce that the subspace $f^{*-1}(\mathfrak{p})$ of $Y$ is naturally homeomorphic to $\Spec(B_{\mathfrak{p}}/\mathfrak{p}B_{\mathfrak{p}}) = \Spec(k(\mathfrak{p}) \otimes_A B)$, where $k(\mathfrak{p})$ is the residue field of the local ring $A_{\mathfrak{p}}$. $\Spec(k(\mathfrak{p}) \otimes_A B)$ is called the \emph{fiber} of $f^*$ over $\mathfrak{p}$.
    \end{enumerate}

    \item \label{exer22-chapter3}Let $A$ be a ring and $\mathfrak{p}$ a prime ideal of $A$. Then the canonical image of $\Spec(A_{\mathfrak{p}})$ in $\Spec(A)$ is equal to the intersection of all the open neighborhoods of $\mathfrak{p}$ in $\Spec(A)$.

    \item \label{exer23-chapter3}Let $A$ be a ring, let $X = \Spec(A)$ and let $U$ be a basic open set in $X$ (i.e., $U = X_f$ for some $f \in A$: Chapter \ref{chapter1}, Exercise \ref{exer17-chapter1}).
    \begin{enumerate}[i)]
        \item If $U = X_f$, show that the ring $A(U) = A_f$ depends only on $U$ and not on $f$.
        \item Let $U' = X_g$ be another basic open set such that $U' \subseteq U$. Show that there is an equation of the form $g^n = uf$ for some integer $n > 0$ and some $u \in A$, and use this to define a homomorphism $\rho \colon A(U) \to A(U')$ (i.e., $A_f \to A_g$) by mapping $a/f^m$ to $au^m/g^{mn}$. Show that $\rho$ depends only on $U$ and $U'$. This homomorphism is called the \emph{restriction homomorphism}.
        \item If $U = U'$, then $\rho$ is the identity map.
        \item If $U \supseteq U' \supseteq U''$ are basic open sets in $X$, show that the diagram (with arrows being restriction homomorphisms) is commutative.
        \[\begin{tikzcd}
            {A(U)} && {A(U'')} \\
            & {A(U')}
            \arrow[from=1-1, to=1-3]
            \arrow[from=1-1, to=2-2]
            \arrow[from=2-2, to=1-3]
        \end{tikzcd}\]
        \item Let $x$ ($= \mathfrak{p}$) be a point of $X$. Show that \[\varinjlim_{U\ni x} A(U) \cong A_{\mathfrak{p}}.\]
    \end{enumerate}
    The assignment of the ring $A(U)$ to each basic open set $U$ of $X$, and the restriction homomorphisms $\rho$, satisfying the conditions iii) and iv) above, constitutes a \emph{presheaf of rings} on the basis of open sets $(X_f)_{f \in A}$. v) says that the \emph{stalk} of this presheaf at $x \in X$ is the corresponding local ring $A_{\mathfrak{p}}$.

    \item \label{exer24-chapter3}Show that the presheaf of Exercise \ref{exer23-chapter3} has the following property. Let $(U_i)_{i \in I}$ be a covering of $X$ by basic open sets. For each $i \in I$ let $s_i \in A(U_i)$ be such that, for each pair of indices $i,j$, the images of $s_i$ and $s_j$ in $A(U_i \cup U_j)$ are equal. Then there exists a unique $s \in A$ ($= A(X)$) whose image in $A(U_i)$ is $s_i$ for all $i \in I$. (This essentially implies that the presheaf is a \emph{sheaf}.)

    \item \label{exer25-chapter3}Let $f \colon A \to B$, $g \colon A \to C$ be ring homomorphisms and let $h \colon A \to B \otimes_A C$ be defined by $h(x) = f(x) \otimes g(x)$. Let $X, Y, Z, T$ be the prime spectra of $A, B, C, B \otimes_A C$ respectively. Then $h^*(T) = f^*(Y) \cap g^*(Z)$.
    
    \textit{Hint:} Let $\mathfrak{p} \in X$, and let $k = k(\mathfrak{p})$ be the residue field at $\mathfrak{p}$. By Exercise 21, the fiber $h^{*-1}(\mathfrak{p})$ is the spectrum of $(B \otimes_A C) \otimes_A k \cong (B \otimes_A k) \otimes_k (C \otimes_A k)$. Hence $\mathfrak{p} \in h^*(T) \iff (B \otimes_A k) \otimes_k (C \otimes_A k) \neq 0 \iff B \otimes_A k \neq 0 \text{ and } C \otimes_A k \neq 0 \iff \mathfrak{p} \in f^*(Y) \cap g^*(Z)$.

    \item \label{exer26-chapter3}Let $(B_\alpha, g_{\alpha\beta})$ be a direct system of rings and $B$ the direct limit. For each $\alpha$, let $f_\alpha \colon A \to B_\alpha$ be a ring homomorphism such that $g_{\alpha\beta} \circ f_\alpha = f_\beta$ whenever $\alpha \leq \beta$ (i.e., the $B_\alpha$ form a direct system of $A$-algebras). The $f_\alpha$ induce $f \colon A \to B$. Show that
    \[
    f^*(\Spec(B)) = \bigcap_\alpha f_\alpha^*(\Spec(B_\alpha)).
    \]
    \textit{Hint:} Let $\mathfrak{p} \in \Spec(A)$. Then $f^{*-1}(\mathfrak{p})$ is the spectrum of \[B \otimes_A k(\mathfrak{p}) \cong \varinjlim (B_\alpha \otimes_A k(\mathfrak{p}))\] (since tensor products commute with direct limits: Chapter \ref{chapter2}, Exercise \ref{exer20-chapter2}). By Exercise \ref{exer21-chapter2} of Chapter \ref{chapter2} it follows that $f^{*-1}(\mathfrak{p}) = \varnothing$ if and only if $B_\alpha \otimes_A k(\mathfrak{p}) = 0$ for some $\alpha$, i.e., if and only if $f_\alpha^{*-1}(\mathfrak{p}) = \varnothing$.

    \item \label{exer27-chapter3} \begin{enumerate}[i)]
        \item Let $f_\alpha \colon A \to B_\alpha$ be any family of $A$-algebras and let $f \colon A \to B$ be their tensor product over $A$ (Chapter \ref{chapter2}, Exercise \ref{exer23-chapter2}). Then \[f^*(\Spec(B)) = \bigcap_\alpha f_\alpha^*(\Spec(B_\alpha)).\] [Use Exercises \ref{exer25-chapter3} and \ref{exer26-chapter3}.]
        \item Let $f_\alpha \colon A \to B_\alpha$ be any finite family of $A$-algebras and let $B = \prod_\alpha B_\alpha$. Define $f \colon A \to B$ by $f(x) = (f_\alpha(x))$. Then $f^*(\Spec(B)) = \bigcup_\alpha f_\alpha^*(\Spec(B_\alpha))$.
        \item Hence the subsets of $X = \Spec(A)$ of the form $f^*(\Spec(B))$, where $f \colon A \to B$ is a ring homomorphism, satisfy the axioms for closed sets in a topological space. The associated topology is the \emph{constructible topology}\index{Constructible topology} on $X$. It is finer than the Zariski topology (i.e., there are more open sets, or equivalently more closed sets).
        \item Let $X_c$ denote the set $X$ endowed with the constructible topology. Show that $X_c$ is quasi-compact.
    \end{enumerate}

    \item \label{exer28-chapter3} (Continuation of Exercise \ref{exer27-chapter3}.)
    \begin{enumerate}[i)]
        \item For each $g \in A$, the set $X_g$ (Chapter \ref{chapter1}, Exercise \ref{exer17-chapter1}) is both open and closed in the constructible topology.
        \item Let $C'$ denote the smallest topology on $X$ for which the sets $X_g$ are both open and closed, and let $X_{C'}$ denote the set $X$ endowed with this topology. Show that $X_{C'}$ is Hausdorff.
        \item Deduce that the identity mapping $X_c \to X_{C'}$ is a homeomorphism. Hence a subset $E$ of $X$ is of the form $f^*(\Spec(B))$ for some $f \colon A \to B$ if and only if it is closed in the topology $C'$.
        \item The topological space $X_c$ is compact, Hausdorff and totally disconnected.
    \end{enumerate}

    \item Let $f \colon A \to B$ be a ring homomorphism. Show that $f^* \colon \Spec(B) \to \Spec(A)$ is a continuous closed mapping (i.e., maps closed sets to closed sets) for the constructible topology.

    \item Show that the Zariski topology and the constructible topology on $\Spec(A)$ are the same if and only if $A/\mathfrak{N}$ is absolutely flat (where $\mathfrak{N}$ is the nilradical of $A$).\\
    \textit{Hint:} Use Exercise \ref{exer11-chapter3}.
\end{enumerate}

